\title{2021 年高考数学 浙江卷}
\date{}
\maketitle

\section*{一、选择题}
本大题共10小题，每小题4分，共40分。在每小题给出的四个选项中，只有一项是符合题目要求的。

\begin{question}[points=4]
$A = \{x \mid x \geq 1\}$, $B = \{x \mid -1 < x < 2\}$, 则 $A \cap B =$ （   ）
\begin{choices}
  \item $\{x \mid x > -1\}$
  \item $\{x \mid x \geq 1\}$
  \item $\{x \mid -1 < x < 1\}$
  \item $\{x \mid 1 \leq x < 2\}$
\end{choices}
\begin{solution}
\textbf{答案：}D

\textbf{分析：}由题意结合交集的定义可得结果.

\textbf{详解：}由交集的定义可得 $A \cap B = \{x \mid 1 \leq x < 2\}$.
\end{solution}
\end{question}

\begin{question}[points=4]
已知 $a \in \mathbb{R}$, $(1+ai)i=3+i$, (i为虚数单位), 则 $a=$ （   ）
\begin{choices}
  \item $-1$
  \item $1$
  \item $-3$
  \item $3$
\end{choices}
\begin{solution}
\textbf{答案：}C

\textbf{分析：}首先计算左侧的结果，然后结合复数相等的充分必要条件即可求得实数 $a$ 的值.

\textbf{详解：}$(1+ai)i = i + ai^2 = i - a = -a + i = 3 + i$, 所以 $-a=3$, 故 $a=-3$.
\end{solution}
\end{question}

\begin{question}[points=4]
已知非零向量 $\vec{a},\vec{b},\vec{c}$, 则 "$\vec{a}\cdot\vec{c}=\vec{b}\cdot\vec{c}$" 是 "$\vec{a}=\vec{b}$" 的（   ）
\begin{choices}
  \item 充分不必要条件
  \item 必要不充分条件
  \item 充分必要条件
  \item 既不充分又不必要条件
\end{choices}
\begin{solution}
\textbf{答案：}B

\textbf{分析：}考虑两者之间的推出关系后可得两者之间的条件关系.

\textbf{详解：}如图, $\overrightarrow{OA}=\vec{a}$, $\overrightarrow{OB}=\vec{b}$, $\overrightarrow{OC}=\vec{c}$, 当 $\vec{a}-\vec{b}$ 与 $\vec{c}$ 垂直时, $\vec{a}\cdot\vec{c}=\vec{b}\cdot\vec{c}$ 但 $\vec{a}\neq\vec{b}$, 故不是充分条件；反之当 $\vec{a}=\vec{b}$ 时, $\vec{a}\cdot\vec{c}=\vec{b}\cdot\vec{c}$ 成立, 故是必要条件. 因此是必要不充分条件.
\end{solution}
\end{question}

\begin{question}[points=4]
某几何体的三视图如图所示，则该几何体的体积是（   ）
\begin{choices}
  \item $\frac{3}{2}$
  \item $3$
  \item $\frac{3\sqrt{2}}{2}$
  \item $3\sqrt{2}$
\end{choices}
\begin{solution}
\textbf{答案：}A

\textbf{分析：}根据三视图可得几何体为四棱柱, 根据棱柱的体积公式可求其体积.

\textbf{详解：}几何体为四棱柱, 高为1, 底面为等腰梯形, 上底为 $\sqrt{2}$, 下底为 $2\sqrt{2}$, 腰长为1, 故梯形的高为 $\sqrt{1^2 - \left(\frac{\sqrt{2}}{2}\right)^2} = \frac{\sqrt{2}}{2}$, 体积 $V = \frac{1}{2}(\sqrt{2}+2\sqrt{2})\times\frac{\sqrt{2}}{2}\times1 = \frac{3}{2}$.
\end{solution}
\end{question}

\begin{question}[points=4]
若实数 $x,y$ 满足约束条件 $\begin{cases} x+1 \geq 0 \\ x-y \leq 0 \\ 2x+3y-1 \leq 0 \end{cases}$, 则 $z = x - \frac{1}{2}y$ 的最小值是（   ）
\begin{choices}
  \item $-2$
  \item $-\frac{3}{2}$
  \item $-\frac{1}{2}$
  \item $\frac{1}{10}$
\end{choices}
\begin{solution}
\textbf{答案：}B

\textbf{分析：}画出可行域, 目标函数化为 $y=2x-2z$, 求过可行域点且斜率为2的直线在 y 轴上截距的最大值.

\textbf{详解：}可行域如图, 由 $\begin{cases} x=-1 \\ 2x+3y-1=0 \end{cases}$ 得 $A(-1,1)$. 当直线 $y=2x-2z$ 过 $A$ 点时, $z$ 取最小值 $z = -1 - \frac{1}{2}\times1 = -\frac{3}{2}$.
\end{solution}
\end{question}

\begin{question}[points=4]
如图已知正方体 $ABCD-A_1B_1C_1D_1$, $M,N$ 分别是 $A_1D$, $D_1B$ 的中点, 则（   ）
\begin{choices}
  \item 直线 $A_1D$ 与直线 $D_1B$ 垂直, 直线 $MN //$ 平面 $ABCD$
  \item 直线 $A_1D$ 与直线 $D_1B$ 平行, 直线 $MN \perp$ 平面 $BDD_1B_1$
  \item 直线 $A_1D$ 与直线 $D_1B$ 相交, 直线 $MN //$ 平面 $ABCD$
  \item 直线 $A_1D$ 与直线 $D_1B$ 异面, 直线 $MN \perp$ 平面 $BDD_1B_1$
\end{choices}
\begin{solution}
\textbf{答案：}A

\textbf{分析：}由正方体间的垂直、平行关系, 可证 $MN // AB$, $A_1D \perp$ 平面 $ABD_1$, 即可得出结论.

\textbf{详解：}连 $AD_1$, $M$ 是 $A_1D$ 中点, 所以 $M$ 为 $AD_1$ 中点, 又 $N$ 是 $D_1B$ 中点, 所以 $MN // AB$, 故 $MN //$ 平面 $ABCD$; 由正方体性质, $A_1D \perp AD_1$, $AB \perp A_1D$, 所以 $A_1D \perp$ 平面 $ABD_1$, 从而 $A_1D \perp D_1B$, 且异面.
\end{solution}
\end{question}

\begin{question}[points=4]
已知函数 $f(x) = x^2 + \frac{1}{4}$, $g(x) = \sin x$, 则图象为如图的函数可能是（   ）
\begin{choices}
  \item $y = f(x) + g(x) - \frac{1}{4}$
  \item $y = f(x) - g(x) - \frac{1}{4}$
  \item $y = f(x)g(x)$
  \item $y = \frac{g(x)}{f(x)}$
\end{choices}
\begin{solution}
\textbf{答案：}D

\textbf{分析：}由函数的奇偶性可排除 A、B, 结合导数判断函数的单调性可判断 C, 即可得解.

\textbf{详解：}A 非奇非偶, 排除；B 非奇非偶, 排除；C 求导在 $x=\frac{\pi}{4}$ 处导数为正, 与图象不符, 排除. 故选 D.
\end{solution}
\end{question}

\begin{question}[points=4]
已知 $\alpha,\beta,\gamma$ 是互不相同的锐角, 则在 $\sin\alpha\cos\beta$, $\sin\beta\cos\gamma$, $\sin\gamma\cos\alpha$ 三个值中, 大于 $\frac{1}{2}$ 的个数的最大值是（   ）
\begin{choices}
  \item $0$
  \item $1$
  \item $2$
  \item $3$
\end{choices}
\begin{solution}
\textbf{答案：}C

\textbf{分析：}利用基本不等式或排序不等式得 $\sin\alpha\cos\beta+\sin\beta\cos\gamma+\sin\gamma\cos\alpha \leq \frac{3}{2}$, 从而可判断三个代数式不可能均大于 $\frac{1}{2}$, 再结合特例可得最大值.

\textbf{详解：}由基本不等式, $\sin\alpha\cos\beta \leq \frac{\sin^2\alpha+\cos^2\beta}{2}$, 同理可得和 $\leq \frac{3}{2}$, 故不可能三个都大于 $\frac{1}{2}$. 取 $\alpha=\frac{\pi}{6},\beta=\frac{\pi}{3},\gamma=\frac{\pi}{4}$, 则 $\sin\alpha\cos\beta = \frac{1}{4}<\frac{1}{2}$, $\sin\beta\cos\gamma = \frac{\sqrt{6}}{4}>\frac{1}{2}$, $\sin\gamma\cos\alpha = \frac{\sqrt{6}}{4}>\frac{1}{2}$, 故最大个数为2.
\end{solution}
\end{question}

\begin{question}[points=4]
已知 $a,b\in\mathbb{R}$, $ab>0$, 函数 $f(x)=ax^2+b(x\in\mathbb{R})$. 若 $f(s-t), f(s), f(s+t)$ 成等比数列, 则平面上点 $(s,t)$ 的轨迹是（   ）
\begin{choices}
  \item 直线和圆
  \item 直线和椭圆
  \item 直线和双曲线
  \item 直线和抛物线
\end{choices}
\begin{solution}
\textbf{答案：}C

\textbf{分析：}利用等比数列建立等式, 然后恒等变形确定轨迹方程.

\textbf{详解：}由 $f(s-t)f(s+t)=[f(s)]^2$ 得 $(a(s-t)^2+b)(a(s+t)^2+b) = (as^2+b)^2$, 整理得 $-2a^2s^2t^2+a^2t^4+2abt^2=0$, 即 $t^2(-2a s^2 + a t^2 + 2b)=0$, 所以 $t=0$ 或 $-2as^2+at^2+2b=0$, 即 $\frac{s^2}{\frac{b}{a}} - \frac{t^2}{\frac{2b}{a}}=1$ (双曲线) 或 $t=0$ (直线).
\end{solution}
\end{question}

\begin{question}[points=4]
已知数列 $\{a_n\}$ 满足 $a_1=1$, $a_{n+1} = \frac{a_n}{1+\sqrt{a_n}}$ $(n\in\mathbb{N}^*)$. 记数列 $\{a_n\}$ 的前 $n$ 项和为 $S_n$, 则（   ）
\begin{choices}
  \item $\frac{1}{2} < S_{100} < 3$
  \item $3 < S_{100} < 4$
  \item $4 < S_{100} < \frac{9}{2}$
  \item $\frac{9}{2} < S_{100} < 5$
\end{choices}
\begin{solution}
\textbf{答案：}A

\textbf{分析：}利用倒数法得到 $\frac{1}{a_{n+1}} = \frac{1}{a_n} + \frac{1}{\sqrt{a_n}}$, 再放缩可得 $a_n$ 的范围, 进而求和.

\textbf{详解：}由 $a_{n+1} = \frac{a_n}{1+\sqrt{a_n}}$ 得 $\frac{1}{a_{n+1}} = \frac{1}{a_n} + \frac{1}{\sqrt{a_n}}$, 得 $a_n \geq \frac{4}{(n+1)^2}$, 且 $a_n \leq \frac{6}{(n+1)(n+2)}$, 求和得 $\frac{1}{2} < S_{100} < 3$.
\end{solution}
\end{question}

\section*{二、填空题}
本大题共7小题，多空题每题6分，单空题每题4分，共36分。

\begin{question}[points=4]
我国古代数学家赵爽用弦图给出了勾股定理的证明. 弦图是由四个全等的直角三角形和中间的一个小正方形拼成的一个大正方形(如图所示). 若直角三角形直角边的长分别是3,4, 记大正方形的面积为 $S_1$, 小正方形的面积为 $S_2$, 则 $\frac{S_1}{S_2} =$ .

\begin{center}
  \includegraphics[width=.38\linewidth]{question-11-1.jpg}
\end{center}
\fillin[{$25$}]
\begin{solution}
\textbf{答案：}$25$

\textbf{分析：}分别求得大正方形的面积和小正方形的面积, 然后计算比值.

\textbf{详解：}大正方形边长 $\sqrt{3^2+4^2}=5$, $S_1=25$; 小正方形面积 $S_2=25-4\times\frac{1}{2}\times3\times4=1$, 故 $\frac{S_1}{S_2}=25$.
\end{solution}
\end{question}

\begin{question}[points=4]
已知 $a\in\mathbb{R}$, 函数 $f(x)=\begin{cases} x^2-4, & x>2 \\ |x-3|+a, & x\leq 2 \end{cases}$, 若 $f[f(\sqrt{6})]=3$, 则 $a=$ .
\fillin[{$2$}]
\begin{solution}
\textbf{答案：}$2$

\textbf{分析：}由题意结合函数的解析式得到关于 $a$ 的方程, 解方程可得 $a$ 的值.

\textbf{详解：}$f(\sqrt{6}) = (\sqrt{6})^2-4=2$, 则 $f[f(\sqrt{6})]=f(2)=|2-3|+a=1+a=3$, 所以 $a=2$.
\end{solution}
\end{question}

\begin{question}[points=6]
已知多项式 $(x-1)^3 + (x+1)^4 = x^4 + a_1 x^3 + a_2 x^2 + a_3 x + a_4$, 则 $a_1 =$ ，$a_2+a_3+a_4 =$ .
\fillin[{$5; 10$}]
\begin{solution}
\textbf{答案：}$5; 10$

\textbf{分析：}根据二项展开式定理, 分别求出 $(x-1)^3$, $(x+1)^4$ 的展开式, 即可得出结论.

\textbf{详解：}$(x-1)^3 = x^3-3x^2+3x-1$, $(x+1)^4 = x^4+4x^3+6x^2+4x+1$, 所以 $a_1=1+4=5$, $a_2=-3+6=3$, $a_3=3+4=7$, $a_4=-1+1=0$, 故 $a_2+a_3+a_4=10$.
\end{solution}
\end{question}

\begin{question}[points=6]
在 $\triangle ABC$ 中, $\angle B = 60^\circ$, $AB=2$, $M$ 是 $BC$ 的中点, $AM=2\sqrt{3}$, 则 $AC=$ ，$\cos\angle MAC=$ .
\fillin[{$2\sqrt{13}$; $\frac{2\sqrt{39}}{13}$}]
\begin{solution}
\textbf{答案：}$2\sqrt{13}$; $\frac{2\sqrt{39}}{13}$

\textbf{分析：}由余弦定理求得 $BC$, 再由余弦定理可得 $AC$ 和 $\cos\angle MAC$.

\textbf{详解：}在 $\triangle ABM$ 中, 由余弦定理: $12=4+BM^2-2\cdot2\cdot BM\cdot\frac12$, 解得 $BM=4$, 则 $BC=8$. 在 $\triangle ABC$ 中, $AC^2=4+64-2\cdot2\cdot8\cdot\frac12=52$, 所以 $AC=2\sqrt{13}$. 在 $\triangle AMC$ 中, $\cos\angle MAC = \frac{AC^2+AM^2-MC^2}{2\cdot AM\cdot AC} = \frac{52+12-16}{2\cdot2\sqrt{3}\cdot2\sqrt{13}} = \frac{48}{8\sqrt{39}} = \frac{2\sqrt{39}}{13}$.
\end{solution}
\end{question}

\begin{question}[points=6]
袋中有4个红球 $m$ 个黄球, $n$ 个绿球. 现从中任取两个球, 记取出的红球数为 $\xi$, 若取出的两个球都是红球的概率为 $\frac{1}{6}$, 一红一黄的概率为 $\frac{1}{3}$, 则 $m-n=$ ，$E(\xi)=$ .
\fillin[{1; $\frac{8}{9}$}]
\begin{solution}
\textbf{答案：}1; $\frac{8}{9}$

\textbf{分析：}根据古典概型的概率公式列式求得 $m,n$, 再求 $\xi$ 的分布列和期望.

\textbf{详解：}$P(\xi=2)=\frac{C_4^2}{C_{m+n+4}^2}=\frac{6}{C_{m+n+4}^2}=\frac16$, 得 $C_{m+n+4}^2=36$, 所以 $m+n+4=9$, 即 $m+n=5$. $P(\text{一红一黄})=\frac{C_4^1 C_m^1}{C_9^2}=\frac{4m}{36}=\frac13$, 解得 $m=3$, 所以 $n=2$, $m-n=1$. $\xi$ 可取0,1,2. $P(\xi=2)=\frac16$, $P(\xi=1)=\frac{C_4^1 C_5^1}{C_9^2}=\frac{20}{36}=\frac59$, $P(\xi=0)=1-\frac16-\frac59=\frac{5}{18}$, $E(\xi)=2\times\frac16+1\times\frac59+0\times\frac{5}{18}=\frac13+\frac59=\frac89$.
\end{solution}
\end{question}

\begin{question}[points=6]
已知椭圆 $\frac{x^2}{a^2}+\frac{y^2}{b^2}=1 (a>b>0)$, 焦点 $F_1(-c,0)$, $F_2(c,0)$ $(c>0)$, 若过 $F_1$ 的直线和圆 $\left(x-\frac12c\right)^2+y^2=c^2$ 相切, 与椭圆在第一象限交于点 $P$, 且 $PF_2 \perp x$ 轴, 则该直线的斜率是，椭圆的离心率是.

\begin{center}
  \includegraphics[width=.38\linewidth]{question-16-1.jpg}
\end{center}
\fillin[{$\frac{2\sqrt{5}}{5}$; $\frac{\sqrt{5}}{5}$}]
\begin{solution}
\textbf{答案：}$\frac{2\sqrt{5}}{5}$; $\frac{\sqrt{5}}{5}$

\textbf{分析：}设 $c=2$, 利用几何关系求得 $\sin\angle PF_1F_2$, 进而得到斜率, 再由椭圆定义求离心率.

\textbf{详解：}如图, 设 $c=2$, 切点为 $B$, $\sin\angle PF_1F_2 = \frac{AB}{F_1A} = \frac{2}{3}$, 则 $\tan\angle PF_1F_2 = \frac{2}{\sqrt{5}}$, 所以斜率 $k=\frac{2\sqrt{5}}{5}$. 又 $F_1F_2=4$, $PF_2 = F_1F_2 \cdot \tan\angle PF_1F_2 = \frac{8\sqrt{5}}{5}$, $PF_1 = \frac{PF_2}{\sin\angle PF_1F_2} = \frac{12\sqrt{5}}{5}$, 所以 $2a = PF_1+PF_2 = 4\sqrt{5}$, $a=2\sqrt{5}$, $e=\frac{c}{a}=\frac{\sqrt{5}}{5}$.
\end{solution}
\end{question}

\begin{question}[points=4]
已知平面向量 $\vec{a},\vec{b},\vec{c}$ $(\vec{c}\neq\vec{0})$ 满足 $|\vec{a}|=1$, $|\vec{b}|=2$, $\vec{a}\cdot\vec{b}=0$, $(\vec{a}-\vec{b})\cdot\vec{c}=0$. 记向量 $\vec{d}$ 在 $\vec{a},\vec{b}$ 方向上的投影分别为 $x,y$, $\vec{d}-\vec{a}$ 在 $\vec{c}$ 方向上的投影为 $z$, 则 $x^2+y^2+z^2$ 的最小值为.
\fillin[{$\frac{2}{5}$}]
\begin{solution}
\textbf{答案：}$\frac{2}{5}$

\textbf{分析：}设 $\vec{a}=(1,0)$, $\vec{b}=(0,2)$, $\vec{c}=(m,n)$, 由条件得 $m=2n$, 再由投影公式得 $2x+y\pm\sqrt{5}z=2$, 利用柯西不等式求最小值.

\textbf{详解：}令 $\vec{a}=(1,0)$, $\vec{b}=(0,2)$, $\vec{c}=(2n,n)$, 设 $\vec{d}=(x,y)$, 则 $z = \frac{(\vec{d}-\vec{a})\cdot\vec{c}}{|\vec{c}|} = \frac{2n(x-1)+ny}{\sqrt{5}n} = \frac{2x+y-2}{\sqrt{5}}$, 故 $2x+y-\sqrt{5}z=2$. 由柯西不等式: $(2^2+1^2+(\sqrt{5})^2)(x^2+y^2+z^2) \geq (2x+y-\sqrt{5}z)^2=4$, 所以 $x^2+y^2+z^2 \geq \frac{4}{10}=\frac{2}{5}$, 等号当 $\frac{x}{2}=\frac{y}{1}=\frac{z}{-\sqrt{5}}$ 时取得, 即 $x=\frac{4}{5},y=\frac{2}{5},z=-\frac{2\sqrt{5}}{5}$.
\end{solution}
\end{question}

\section*{三、解答题}
本大题共5小题，共74分。解答应写出文字说明、证明过程或演算步骤。

\begin{question}[points=14]
设函数 $f(x)=\sin x+\cos x$ $(x\in\mathbb{R})$.

(1) 求函数 $y=\left[f\left(x+\frac{\pi}{2}\right)\right]^2$ 的最小正周期；

(2) 求函数 $y=f(x)f\left(x-\frac{\pi}{4}\right)$ 在 $[0,\frac{\pi}{2}]$ 上的最大值.
\begin{solution}
\textbf{答案：}(1) $\pi$; (2) $1+\frac{\sqrt{2}}{2}$

\textbf{分析：}(1) 利用三角恒等变换化简, 再求周期; (2) 化简表达式, 利用正弦函数的性质求最大值.

\textbf{详解：}(1) $f(x)=\sqrt{2}\sin\left(x+\frac{\pi}{4}\right)$, 则 $y=\left[\sqrt{2}\sin\left(x+\frac{3\pi}{4}\right)\right]^2=2\sin^2\left(x+\frac{3\pi}{4}\right)=1-\cos(2x+\frac{3\pi}{2})=1-\sin2x$, 故最小正周期 $T=\frac{2\pi}{2}=\pi$.
(2) $y=f(x)f(x-\frac{\pi}{4})=\sqrt{2}\sin(x+\frac{\pi}{4})\cdot\sqrt{2}\sin x = 2\sin x\cdot\frac{\sqrt{2}}{2}(\sin x+\cos x) = \sqrt{2}\sin^2 x + \sqrt{2}\sin x\cos x = \frac{\sqrt{2}}{2}(1-\cos2x) + \frac{\sqrt{2}}{2}\sin2x = \sin(2x-\frac{\pi}{4})+\frac{\sqrt{2}}{2}$. 当 $x\in[0,\frac{\pi}{2}]$ 时, $2x-\frac{\pi}{4}\in[-\frac{\pi}{4},\frac{3\pi}{4}]$, 故当 $2x-\frac{\pi}{4}=\frac{\pi}{2}$ 即 $x=\frac{3\pi}{8}$ 时, $y$ 取最大值 $1+\frac{\sqrt{2}}{2}$.
\end{solution}
\end{question}

\begin{question}[points=15]
如图，在四棱锥 $P-ABCD$ 中，底面 $ABCD$ 是平行四边形，$\angle ABC=120^\circ$, $AB=1$, $BC=4$, $PA=\sqrt{15}$, $M,N$ 分别为 $BC, PC$ 的中点，$PD\perp DC$, $PM\perp MD$.

(1) 证明：$AB\perp PM$；

(2) 求直线 $AN$ 与平面 $PDM$ 所成角的正弦值.
\begin{solution}
\textbf{答案：}(1) 见解析; (2) $\frac{\sqrt{15}}{6}$

\textbf{分析：}(1) 先证 $DC\perp$ 平面 $PDM$, 再由 $AB\parallel DC$ 得 $AB\perp PM$; (2) 建立空间直角坐标系, 利用向量法求线面角.

\textbf{详解：}(1) 在 $\triangle DCM$ 中, $DC=1,CM=2,\angle DCM=60^\circ$, 由余弦定理得 $DM=\sqrt{3}$. 由 $DM^2+DC^2=CM^2$ 得 $DM\perp DC$. 又 $PD\perp DC$, $PD\cap DM=D$, 故 $DC\perp$ 平面 $PDM$, 从而 $DC\perp PM$. 又 $AB\parallel DC$, 所以 $AB\perp PM$.
(2) 由 $PM\perp MD$ 及 (1) 得 $PM\perp$ 平面 $ABCD$. 以 $M$ 为原点, $MD$ 为 $x$ 轴, $ME$ ( $E$ 为 $AD$ 中点) 为 $y$ 轴, $MP$ 为 $z$ 轴建立空间直角坐标系. 则 $A(-\sqrt{3},2,0)$, $D(\sqrt{3},0,0)$, $P(0,0,2\sqrt{2})$, $C(\sqrt{3},-1,0)$, $N(\frac{\sqrt{3}}{2},-\frac12,\sqrt{2})$, 所以 $\overrightarrow{AN}=(\frac{3\sqrt{3}}{2},-\frac52,\sqrt{2})$. 平面 $PDM$ 的法向量为 $\vec{n}=(0,1,0)$. 则 $\sin\theta = \frac{|\overrightarrow{AN}\cdot\vec{n}|}{|\overrightarrow{AN}|\cdot|\vec{n}|} = \frac{5/2}{\sqrt{\frac{27}{4}+\frac{25}{4}+2}} = \frac{5/2}{\sqrt{15}} = \frac{\sqrt{15}}{6}$.
\end{solution}
\end{question}

\begin{question}[points=15]
已知数列 $\{a_n\}$ 的前 $n$ 项和为 $S_n$, $a_1=-\frac94$, 且 $4S_{n+1}=3S_n-9$.

(1) 求数列 $\{a_n\}$ 的通项；

(2) 设数列 $\{b_n\}$ 满足 $3b_n+(n-4)a_n=0$ $(n\in\mathbb{N}^*)$, 记 $\{b_n\}$ 的前 $n$ 项和为 $T_n$, 若 $T_n \leq \lambda b_n$ 对任意 $n\in\mathbb{N}^*$ 恒成立, 求实数 $\lambda$ 的取值范围.
\begin{solution}
\textbf{答案：}(1) $a_n = -3\cdot\left(\frac34\right)^n$; (2) $[-3,1]$

\textbf{分析：}(1) 利用 $S_n$ 与 $a_n$ 的关系, 得到等比数列; (2) 先求 $b_n$, 再用错位相减法求 $T_n$, 再分离参数 $\lambda$ 求解.

\textbf{详解：}(1) 当 $n=1$ 时, $4(a_1+a_2)=3a_1-9$, 得 $a_2=-\frac{27}{16}$. 当 $n\geq2$ 时, $4S_{n+1}=3S_n-9$ 与 $4S_n=3S_{n-1}-9$ 相减得 $4a_{n+1}=3a_n$, 又 $\frac{a_2}{a_1}=\frac34$, 故 $\{a_n\}$ 是首项为 $-\frac94$, 公比为 $\frac34$ 的等比数列, 所以 $a_n = -\frac94\cdot\left(\frac34\right)^{n-1} = -3\cdot\left(\frac34\right)^n$.
(2) 由 $3b_n+(n-4)a_n=0$ 得 $b_n = (n-4)\left(\frac34\right)^n$. 则 $T_n = -3\cdot\frac34 -2\cdot\left(\frac34\right)^2 -1\cdot\left(\frac34\right)^3 +0\cdot\left(\frac34\right)^4 +\cdots+(n-4)\left(\frac34\right)^n$, 错位相减得 $T_n = -4n\left(\frac34\right)^{n+1}$. 由 $T_n\leq\lambda b_n$ 得 $-4n\left(\frac34\right)^{n+1} \leq \lambda (n-4)\left(\frac34\right)^n$, 即 $\lambda(n-4) + 3n \geq 0$ 恒成立. 当 $n=4$ 时显然成立；当 $n<4$ 时, $\lambda \leq -\frac{3n}{n-4} = -3-\frac{12}{n-4}$, 右侧最大值为1 (当 $n=3$), 故 $\lambda\leq1$；当 $n>4$ 时, $\lambda \geq -\frac{3n}{n-4} = -3-\frac{12}{n-4}$, 右侧最小值为-3 (当 $n\to\infty$), 故 $\lambda\geq-3$. 综上, $-3\leq\lambda\leq1$.
\end{solution}
\end{question}

\begin{question}[points=15]
如图，已知 $F$ 是抛物线 $y^2=2px (p>0)$ 的焦点，$M$ 是抛物线的准线与 $x$ 轴的交点，且 $MF=2$.

(1) 求抛物线的方程；

(2) 设过点 $F$ 的直线交抛物线于 $A,B$ 两点, 斜率为 $2$ 的直线 $l$ 与直线 $MA, MB, AB$, $x$ 轴依次交于点 $P,Q,R,N$, 且 $RN^2 = PN\cdot QN$, 求直线 $l$ 在 $x$ 轴上截距的范围.
\begin{solution}
\textbf{答案：}(1) $y^2=4x$; (2) $(-\infty, -7-4\sqrt{3}] \cup [-7+4\sqrt{3},1) \cup (1,+\infty)$

\textbf{分析：}(1) 由 $MF=2$ 得 $p=2$; (2) 联立直线 $AB$ 与抛物线, 设 $N(n,0)$, 利用 $RN^2=PN\cdot QN$ 转化为坐标关系, 得到 $\left(\frac{n+1}{n-1}\right)^2 = \frac{3+4t^2}{(2t-1)^2}$, 再求范围.

\textbf{详解：}(1) $MF = p = 2$, 故抛物线方程为 $y^2=4x$.
(2) 设 $AB: x=ty+1$, $A(x_1,y_1), B(x_2,y_2)$, $N(n,0)$. 由 $\begin{cases}x=ty+1\\ y^2=4x\end{cases}$ 得 $y^2-4ty-4=0$, $y_1+y_2=4t$, $y_1y_2=-4$. 直线 $MA: y=\frac{y_1}{x_1+1}(x+1)$, 联立 $l: x=\frac{y}{2}+n$ 得 $y_P = \frac{2(n+1)y_1}{2x_1+2-y_1}$, 同理 $y_Q = \frac{2(n+1)y_2}{2x_2+2-y_2}$. 联立 $AB$ 与 $l$ 得 $y_R = \frac{2(n-1)}{2t-1}$. 由 $RN^2=PN\cdot QN$ 得 $y_R^2 = y_P y_Q$, 代入整理得 $\left(\frac{n+1}{n-1}\right)^2 = \frac{3+4t^2}{(2t-1)^2}$. 令 $s=2t-1$, 则 $\frac{3+4t^2}{(2t-1)^2} = \frac{s^2+2s+4}{s^2} = 1+\frac{2}{s}+\frac{4}{s^2} = 4\left(\frac{1}{s}+\frac14\right)^2 + \frac34 \geq \frac34$. 故 $\left(\frac{n+1}{n-1}\right)^2 \geq \frac34$, 且 $n\neq1$, 解得 $n\leq -7-4\sqrt{3}$ 或 $n\geq -7+4\sqrt{3}$ 且 $n\neq1$. 又 $n=1$ 时 $y_R=0$, 此时 $RN^2=0$, $PN\cdot QN=0$ 成立? 实际 $n=1$ 时直线 $l$ 过 $F$, 但由过程 $2t-1\neq0$, 故 $n=1$ 排除. 所以截距范围为 $(-\infty, -7-4\sqrt{3}] \cup [-7+4\sqrt{3},1) \cup (1,+\infty)$.
\end{solution}
\end{question}

\begin{question}[points=15]
设 $a,b$ 为实数, 且 $a>1$, 函数 $f(x)=a^x - bx + e^2$ $(x\in\mathbb{R})$.

(1) 求函数 $f(x)$ 的单调区间；

(2) 若对任意 $b>2e^2$, 函数 $f(x)$ 有两个不同的零点, 求 $a$ 的取值范围；

(3) 当 $a=e$ 时, 证明：对任意 $b>e^4$, 函数 $f(x)$ 有两个不同的零点 $x_1,x_2$, 满足 $x_2 > \frac{b\ln b}{2e^2} x_1 + \frac{e^2}{b}$.

(注：$e=2.71828\cdots$ 是自然对数的底数)
\begin{solution}
\textbf{答案：}(1) 见解析; (2) $(1,e^2]$; (3) 见解析

\textbf{分析：}(1) 求导, 讨论 $b$ 的正负; (2) 转化为方程有解, 利用导数研究函数; (3) 利用 (2) 的结论, 结合放缩法证明.

\textbf{详解：}(1) $f'(x)=a^x\ln a - b$. 若 $b\leq0$, 则 $f'(x)>0$, $f(x)$ 在 $\mathbb{R}$ 上单调递增；若 $b>0$, 则当 $x<\log_a\frac{b}{\ln a}$ 时 $f'(x)<0$, $f(x)$ 单调递减；当 $x>\log_a\frac{b}{\ln a}$ 时 $f'(x)>0$, $f(x)$ 单调递增.
(2) $f(x)=0$ 有两个不同零点等价于 $e^{x\ln a} - bx + e^2=0$ 有两个不同解. 令 $t=x\ln a$, 则 $e^t - \frac{b}{\ln a}t + e^2=0$, 即 $\frac{b}{\ln a} = \frac{e^t+e^2}{t}$. 设 $g(t)=\frac{e^t+e^2}{t}$, $g'(t)=\frac{e^t(t-1)-e^2}{t^2}$, 记 $h(t)=e^t(t-1)-e^2$, $h'(t)=e^t t>0$, 且 $h(2)=0$, 故 $g(t)$ 在 $(0,2)$ 递减, $(2,+\infty)$ 递增, $g(2)=e^2$. 所以 $\frac{b}{\ln a} > e^2$, 即 $\ln a < \frac{b}{e^2}$. 对任意 $b>2e^2$, 需 $\frac{b}{e^2}>2$, 则 $\ln a \leq 2$, 故 $1<a\leq e^2$.
(3) $a=e$ 时, $f(x)=e^x - bx + e^2$. 由 (2) 知有两正根 $x_1<x_2$, 且 $b=\frac{e^{x_1}+e^2}{x_1} = \frac{e^{x_2}+e^2}{x_2} > e^4$. 由 $b>e^4$ 得 $x_2>5$. 由 $b=\frac{e^{x_1}+e^2}{x_1} < \frac{2e^2}{x_1}$ (因为 $e^{x_1}<e^2$ 故不严谨, 实际上 $x_1<2$, 所以 $e^{x_1}<e^2$ 不成立, 需另证). 由 $b=\frac{e^{x_1}+e^2}{x_1} > \frac{e^2}{x_1}$ 得 $x_1 > \frac{e^2}{b}$, 但我们需要 $x_1<\frac{2e^2}{b}$? 事实上, 由 $b=\frac{e^{x_2}+e^2}{x_2} > \frac{e^{x_2}}{x_2}$ 可得 $\ln b > x_2 - \ln x_2$, 即 $x_2 < \ln b + \ln x_2$. 要证 $x_2 > \frac{b\ln b}{2e^2}x_1 + \frac{e^2}{b}$, 先由 $x_1 < \frac{2e^2}{b}$? 实际上由 $b=\frac{e^{x_1}+e^2}{x_1}$ 且 $x_1<2$, $e^{x_1}<e^2$, 得 $b>\frac{2e^2}{x_1}$, 故 $x_1<\frac{2e^2}{b}$. 则右边 $< \frac{b\ln b}{2e^2}\cdot\frac{2e^2}{b}+\frac{e^2}{b} = \ln b + \frac{e^2}{b}$. 而 $b = \frac{e^{x_2}+e^2}{x_2} < \frac{2e^{x_2}}{x_2}$ (因为 $e^2<e^{x_2}$ 当 $x_2>2$), 故 $\ln b < x_2 + \ln2 - \ln x_2$, 只需证 $x_2 > x_2 + \ln2 - \ln x_2 + \frac{e^2}{b}$, 即 $\ln x_2 > \ln2 + \frac{e^2}{b}$, 当 $b>e^4$ 时 $\frac{e^2}{b}<\frac{1}{e^2}$, $\ln x_2 > \ln5 > \ln2 + \frac{1}{e^2}$ 成立. 故原不等式成立.
\end{solution}
\end{question}

